\section{Contributions}
\label{sec:conclusion_contributions}

The first contribution of this study is the formulation of a learning-progress intrinsic-reward family that combines latent transition impact with local predictive improvement.
The method does not reward prediction error directly.
Instead, it rewards transitions in which the learned representation changes and the corresponding latent region shows evidence of decreasing forward-model error.
This formulation addresses the limitation of curiosity signals that treat surprise as intrinsically valuable even when additional experience does not lead to improved prediction.

The second contribution is an online latent-space partitioning procedure for estimating local learning progress during on-policy reinforcement learning.
The partition stores region-specific prediction-error statistics and increases resolution in visited portions of the latent space through bounded recursive splitting.
This allows the method to distinguish heterogeneous dynamics without requiring a fixed hand-designed state abstraction.
The use of short- and long-horizon exponential moving averages provides a compact estimate of whether predictive competence is improving within each region.

The third contribution is a region-wise gating mechanism for suppressing unproductive intrinsic reward.
The gate compares each region's learning progress and recent prediction error against robust references computed over visited regions, and it suppresses the base intrinsic score when a region remains high-error and low-progress for a sufficient duration.
This mechanism operationalizes the distinction between learnable difficulty and unproductive unpredictability.
The empirical results show that this guardrail can improve sparse-reward reliability while also revealing the conditions under which it may become overly conservative.

The fourth contribution is a PPO-compatible integration of the proposed intrinsic reward with explicit reward-scale controls.
The method uses component-wise normalization, intrinsic clipping, terminal-transition suppression, and a cosine taper for the intrinsic coefficient.
These design choices make the intrinsic reward usable as a transient training aid rather than as a permanent replacement for the task objective.
The resulting procedure can be compared directly with established intrinsic-reward baselines under the same policy-optimization backbone.

The fifth contribution is an empirical evaluation that considers performance, reliability, ablation behavior, gate dynamics, and computational efficiency together.
The benchmark suite covers both sparse-reward and continuous-control regimes, and the evaluation reports deterministic extrinsic return from offline checkpoints rather than training reward.
The inclusion of step-AUC, final-checkpoint return, thresholded reliability, wall-clock AUC, and timing breakdowns provides a broad assessment of how intrinsic rewards affect both learning quality and training cost.

The sixth contribution is the characterization of when the proposed components are useful.
The results show that the combined impact--learning-progress reward is most beneficial in sparse-reward exploration, that single components can be competitive in some dense-control settings, and that gating is advantageous only when suppression of unproductive regions outweighs its conservatism and overhead.
This characterization is important because it frames the method as an environment-dependent exploration mechanism rather than as a universally dominant algorithm.